\chapter{METODOLOGI}

\section{Metodologi Penelitian}

\noindent Bab ini menguraikan tahapan metodologi, kerangka desain, dan arsitektur sistem untuk pengembangan model \textit{Credit Risk Scoring} terautomasi dengan integrasi prinsip \textit{Machine Learning Operations} (MLOps)[cite: 3]. Fokus utama sistem ini adalah mengatasi tantangan \textit{data drift} dan \textit{concept drift} secara \textit{real-time}[cite: 4]. Secara umum, metodologi penelitian dirancang untuk memastikan bahwa proses pengembangan model tidak berhenti pada tahap eksperimen, tetapi berlanjut hingga implementasi, pemantauan, dan perbaikan model secara berkesinambungan. Dengan demikian, penelitian ini tidak hanya menghasilkan model prediktif, namun juga menghasilkan kerangka kerja operasional yang siap diterapkan pada lingkungan produksi.

Pendekatan yang digunakan adalah \textit{applied research}, yaitu penelitian terapan yang berorientasi pada penyelesaian masalah nyata melalui perancangan, pembangunan, dan evaluasi sistem[cite: 7, 8]. Permasalahan utama yang diangkat adalah kebutuhan akan sistem penilaian risiko kredit yang akurat, dapat dijelaskan, dan berkelanjutan. Ketiga aspek tersebut menjadi krusial ketika distribusi data berubah seiring waktu. Oleh karena itu, metodologi penelitian ini mengintegrasikan rekayasa data, pengembangan model, evaluasi kualitas data sintetis, serta otomasi proses MLOps dalam satu alur kerja terpadu.


\subsection{Tahapan Penelitian}
Tahapan penelitian disusun secara berurutan agar setiap proses memiliki keluaran yang menjadi masukan bagi tahap berikutnya. Secara garis besar, penelitian ini terdiri atas enam tahap utama, yaitu: (1) perumusan kebutuhan sistem dan identifikasi variabel penelitian, (2) pengumpulan serta simulasi data, (3) pra-pemrosesan dan rekayasa fitur, (4) pelatihan dan evaluasi model, (5) validasi kualitas data sintetis, dan (6) implementasi \textit{pipeline} MLOps untuk \textit{deployment}, \textit{monitoring}, dan \textit{retraining} otomatis. Urutan tahapan ini dipilih agar kualitas model dan kesiapan operasional dievaluasi sejak awal sebelum sistem diimplementasikan secara penuh.

Pada tahap pertama, kebutuhan sistem diidentifikasi berdasarkan karakteristik umum proses analisis risiko kredit pada lembaga keuangan. Pada tahap ini, tujuan model ditetapkan, target prediksi didefinisikan, kendala regulasi diidentifikasi, dan kebutuhan interpretabilitas model ditentukan. Tahap kedua berfokus pada penyusunan dataset sintetis yang merepresentasikan profil nasabah. Tahap ketiga mencakup transformasi data mentah menjadi fitur yang siap dipelajari model. Tahap keempat memusatkan perhatian pada pelatihan model \textit{machine learning}, optimasi parameter, dan evaluasi performa. Tahap kelima berfungsi sebagai \textit{quality gate} untuk memastikan data sintetis memiliki utilitas dan keamanan yang memadai. Tahap keenam mengintegrasikan seluruh komponen ke dalam arsitektur MLOps sehingga model dapat dioperasikan secara berkelanjutan.

\subsection{Desain Sistem dan Arsitektur Penelitian}
Arsitektur sistem dirancang menggunakan pendekatan \textit{pipeline machine learning end-to-end} yang bersifat modular[cite: 9]. Modularitas ini memungkinkan setiap komponen dikembangkan, diuji, dan dipelihara secara terpisah. Pemisahan tersebut mendukung skalabilitas sistem, mempermudah pelacakan eksperimen, serta mengurangi ketergantungan antarproses ketika terjadi pembaruan model atau perubahan data.

Secara konseptual, arsitektur penelitian terdiri atas empat lapisan utama. Lapisan pertama adalah \textit{data layer} yang menangani pembentukan data sintetis, validasi struktur data, dan penyimpanan dataset kerja. Lapisan kedua adalah \textit{modeling layer} yang mencakup transformasi fitur, pelatihan model, evaluasi kinerja, dan penyimpanan artefak model. Lapisan ketiga adalah \textit{quality assurance layer} yang memastikan bahwa dataset sintetis memenuhi kriteria utilitas dan privasi sebelum digunakan lebih lanjut. Lapisan keempat adalah \textit{MLOps layer} yang bertugas mengelola orkestrasi \textit{workflow}, registrasi model, \textit{deployment} layanan inferensi, pemantauan \textit{drift}, dan \textit{retraining} otomatis.

\begin{table}[H]
\centering
\caption{Komponen Arsitektur dan Teknologi}
\label{tab:arsitektur}
\begin{tabularx}{\textwidth}{|p{3.2cm}|X|p{4cm}|}
\hline
\textbf{Komponen} & \textbf{Deskripsi} & \textbf{Teknologi} \\ \hline
\textit{Data Engineering} & Transformasi fitur, validasi struktur data, dan penanganan ketidakseimbangan kelas melalui pendekatan statistik dan teknik oversampling[cite: 14]. & Python, Pandas, Scikit-learn [cite: 14] \\ \hline
\textit{Model Development} & Pelatihan algoritma XGBoost, optimasi \textit{hyperparameter}, dan evaluasi performa model klasifikasi[cite: 15]. & XGBoost, RandomizedSearchCV [cite: 15] \\ \hline
\textit{Quality Gate} & Validasi utilitas data sintetis, evaluasi kesamaan distribusi, dan pengujian privasi untuk mencegah \textit{memorization}[cite: 15]. & Scipy, Scikit-learn, KNN [cite: 15] \\ \hline
\textit{MLOps Automation} & Orkestrasi \textit{workflow}, pelacakan eksperimen, registrasi model, \textit{deployment} API, serta \textit{monitoring} dan \textit{retraining} otomatis[cite: 15]. & n8n, MLflow, FastAPI, Docker [cite: 15] \\ \hline
\end{tabularx}
\end{table}

Dalam implementasinya, n8n berperan sebagai orkestrator proses yang menghubungkan tahapan data, model, dan \textit{deployment}. MLflow digunakan untuk mencatat parameter eksperimen, metrik evaluasi, serta versi model terbaik. FastAPI digunakan untuk menyediakan endpoint inferensi agar model dapat diakses sebagai layanan, sedangkan Docker memastikan lingkungan eksekusi bersifat konsisten dan mudah dipindahkan ke server lain. Integrasi seluruh komponen tersebut mendukung pembentukan \textit{pipeline} yang terdokumentasi, dapat direproduksi, dan siap digunakan pada skenario produksi.


\subsection{Perancangan Database dan Penyimpanan Data}

Sistem menyimpan seluruh data nasabah secara persisten menggunakan \textbf{SQLite} dengan nama file nasabah.db yang disimpan di direktori serving/. Database diinisialisasi secara otomatis pada saat \textit{startup} layanan oleh fungsi init\_db() melalui modul database.py, yang mengimpor seluruh baris dari dataset sintetis synthetic\_credit\_risk\_data.csv apabila tabel belum tersedia.

Database terdiri atas dua tabel. Tabel nasabah merupakan tabel utama yang menyimpan 40 lebih atribut setiap nasabah, mulai dari data identitas, profil pekerjaan, kondisi finansial, parameter risiko hasil model. Tabel sent\_emails berfungsi sebagai log audit yang mencatat setiap pengiriman email konfirmasi kepada pengambil keputusan. Relasi antara kedua tabel ini diikat melalui kolom customer\_id sebagai \textit{foreign key}.

\begin{table}[H]
\centering
\caption{Kelompok Kolom Utama Tabel nasabah}
\label{tab:db-cols}
\begin{tabularx}{\textwidth}{|p{3.5cm}|X|}
\hline
\textbf{Kelompok} & \textbf{Kolom Utama} \\ \hline
Identitas & customer\_id (PK), loan\_id, name\_masked, age, gender, marital\_status \\ \hline
Pekerjaan \& Pendapatan & employment\_type, pension\_status, monthly\_income, salary\_source, income\_stability, salary\_volatility \\ \hline
Kesehatan Finansial & bureau\_score, DTI, past\_due\_months, credit\_history\_length, avg\_account\_balance, other\_active\_loans\_count \\ \hline
Detail Pinjaman & product\_type, loan\_amount, term\_months, interest\_rate, collateral\_flag, ltv, outstanding\_balance \\ \hline
Output Model ML & PD, LGD, EAD, expected\_loss, risk\_grade, missed\_payment\_flag \\ \hline
Status Keputusan & approval\_status (Pending / Accepted / Rejected) \\ \hline
\end{tabularx}
\end{table}

Seluruh operasi \textit{Create, Read, Update, Delete} (CRUD) terhadap tabel nasabah diekspos melalui REST API berbasis \textbf{FastAPI} (app.py). Tabel~\ref{tab:crud-bab3} merangkum pemetaan \textit{endpoint} API ke fungsi database.

\begin{table}[H]
\centering
\caption{Mapping \textit{Endpoint} REST API ke Fungsi Database}
\label{tab:crud-bab3}
\begin{tabularx}{\textwidth}{|l|X|l|}
\hline
\textbf{Metode} & \textbf{Endpoint} & \textbf{Fungsi Database} \\ \hline
GET    & /api/nasabah                          & get\_nasabah\_list() \\ \hline
GET    & /api/nasabah/\{id\}                   & get\_nasabah\_by\_id() \\ \hline
POST   & /api/nasabah                          & create\_nasabah() \\ \hline
PUT    & /api/nasabah/\{id\}                   & update\_nasabah() \\ \hline
DELETE & /api/nasabah/\{id\}                   & delete\_nasabah() \\ \hline
GET    & /api/nasabah/\{id\/\allowbreak approve}           & update\_approval\_status("Accepted") \\ \hline
GET    & /api/nasabah/\{id\/\allowbreak reject}            & update\_approval\_status("Rejected") \\ \hline
POST   & /api/nasabah/\{id\/\allowbreak send-confirmation} & create\_sent\_email() \\ \hline
GET    & /api/sent-emails                      & get\_sent\_emails() \\ \hline
\end{tabularx}
\end{table}

Agar API tahan terhadap ketidakcocokan skema, sebelum setiap operasi INSERT atau UPDATE, modul database.py mengkueri kolom yang tersedia melalui PRAGMA table\_info(nasabah) sehingga hanya kolom yang benar-benar ada di skema yang ditulis. Kolom approval\_status yang tidak terdapat pada dataset CSV asli ditambahkan secara dinamis melalui ALTER TABLE saat init\_db() pertama kali dijalankan.

\subsection{Pengumpulan dan Simulasi Data}
Penelitian ini menggunakan \textit{Synthetic Data Generation} untuk mengatasi keterbatasan akses terhadap data kredit riil yang dilindungi oleh regulasi perlindungan data pribadi[cite: 24]. Data yang dibangkitkan terdiri atas 10.000 sampel yang merefleksikan karakteristik nasabah perbankan ritel, khususnya segmen ASN dan pensiunan[cite: 25, 27]. Pemilihan data sintetis dilakukan agar penelitian tetap dapat merepresentasikan pola umum kredit konsumtif tanpa melanggar aspek kerahasiaan dan kepatuhan hukum.

Pemilihan variabel dalam dataset sintetis didasarkan pada relevansinya terhadap proses penilaian risiko kredit secara praktis. Variabel yang disimulasikan mencakup atribut demografis, atribut finansial, dan atribut perilaku pembayaran. Atribut demografis meliputi usia, status pekerjaan, dan masa kerja. Atribut finansial meliputi pendapatan bulanan, rasio angsuran terhadap pendapatan, jumlah pinjaman, dan tenor pinjaman. Atribut perilaku mencakup riwayat keterlambatan pembayaran dan intensitas penggunaan fasilitas kredit. Struktur variabel tersebut memiliki keterkaitan langsung dengan kemungkinan gagal bayar.

\subsubsection{Dasar Perancangan Struktur Dataset Sintetis}
Meskipun dataset yang digunakan bersifat sintetis, struktur variabel dan skema atribut tidak disusun secara acak sepenuhnya. Penentuan kolom, tipe data, format penulisan, serta hubungan antaratribut dilakukan berdasarkan arahan dan masukan praktis dari pihak industri, khususnya selama pelaksanaan kegiatan magang di Bank Mandiri Taspen. Struktur atribut disesuaikan dengan karakteristik umum data nasabah kredit konsumtif, terutama pada segmen ASN dan pensiunan.

Daftar variabel yang digunakan mencakup atribut identitas nasabah, profil pekerjaan, kondisi finansial, histori kredit, parameter risiko, serta indikator makroekonomi yang relevan terhadap proses analisis risiko kredit. Beberapa atribut utama meliputi pendapatan bulanan, \textit{debt-to-income ratio} (DTI), \textit{probability of default} (PD), \textit{exposure at default} (EAD), \textit{loss given default} (LGD), status pensiun, jenis pekerjaan, skor kredit biro, histori tunggakan, dan informasi produk kredit. Dengan rancangan tersebut, dataset sintetis diharapkan mampu merepresentasikan proses penilaian risiko kredit secara lebih realistis.

Format dan tipe data setiap kolom juga disesuaikan dengan praktik umum pencatatan data pada sistem perbankan. Atribut numerik seperti pendapatan dan saldo rekening direpresentasikan dalam format mata uang Rupiah, atribut kategorikal seperti status pernikahan dan jenis pekerjaan menggunakan label kategoris, sedangkan atribut risiko seperti PD dan LGD direpresentasikan dalam bentuk persentase. Penyesuaian ini dilakukan agar dataset sintetis memiliki struktur yang mendekati implementasi nyata pada lingkungan operasional lembaga keuangan.

\subsubsection{Prosedur Pembentukan dan Validasi Awal Dataset}
Proses pengisian nilai data sintetis dilakukan melalui kombinasi pendekatan statistik dan studi literatur dari berbagai sumber terbuka, seperti laporan industri keuangan, publikasi Otoritas Jasa Keuangan (OJK), dokumentasi \textit{credit scoring}, dataset publik terkait risiko kredit, serta referensi penelitian terdahulu mengenai \textit{credit risk modeling}. Distribusi nilai pada beberapa variabel utama disesuaikan agar tetap berada dalam rentang yang realistis berdasarkan karakteristik umum nasabah kredit ritel di Indonesia.

Sebagai validasi awal, sebanyak 50 baris sampel data sintetis disusun dan diberikan kepada pembimbing lapangan atau atasan selama kegiatan magang di Bank Mandiri Taspen untuk memperoleh penilaian terhadap struktur dan kewajaran data. Verifikasi tersebut bertujuan memastikan bahwa format atribut, konsistensi relasi antarkolom, serta pola data yang dihasilkan masih sesuai dengan karakteristik umum data kredit pada praktik perbankan. Pendekatan ini dilakukan agar penelitian tetap menjaga kerahasiaan data nasabah dan kepatuhan terhadap regulasi perlindungan data pribadi, tanpa menghilangkan relevansi praktis terhadap implementasi sistem \textit{credit risk scoring}.

\subsubsection{Formulasi Probabilitas Default (PD)}
Penentuan label \textit{missed payment} dilakukan menggunakan fungsi probabilitas dasar melalui model logistik laten[cite: 32, 33]. Fungsi ini memetakan kombinasi atribut tertentu, seperti pendapatan, usia, dan rasio kewajiban, ke dalam probabilitas gagal bayar. Secara matematis, formulasi probabilitas default dinyatakan sebagai berikut:
\begin{equation}
    PD = \frac{1}{1 + e^{-f(\text{Income, Age})}}
\end{equation}
[cite: 34, 37]

Fungsi laten pada persamaan tersebut merepresentasikan hubungan nonlinier antara karakteristik nasabah dan risiko kredit. Dalam simulasi ini, label kelas minoritas ditetapkan sebesar 14.8\% untuk merepresentasikan kondisi portofolio kredit yang tidak seimbang namun masih realistis[cite: 38]. Penetapan proporsi ini penting karena model klasifikasi kredit pada praktik nyata umumnya menghadapi jumlah kasus gagal bayar yang jauh lebih kecil dibandingkan kasus lancar bayar. Oleh sebab itu, struktur data yang dibangun perlu mencerminkan tantangan tersebut agar evaluasi model menjadi lebih relevan.

\subsubsection{Pembagian Dataset}
Setelah data sintetis dibentuk, dataset dibagi menjadi data latih, data validasi, dan data uji. Pembagian ini dilakukan untuk memisahkan fungsi pelatihan model, pemilihan parameter, dan evaluasi akhir sehingga tidak terjadi kebocoran informasi antarproses. Data latih digunakan untuk membentuk pola prediksi, data validasi digunakan untuk membandingkan konfigurasi model, sedangkan data uji digunakan sebagai dasar penilaian performa akhir model yang bersifat lebih objektif. Dengan mekanisme ini, hasil evaluasi yang diperoleh diharapkan lebih representatif terhadap kondisi implementasi sebenarnya.

\subsection{Pra-pemrosesan Data}
Pra-pemrosesan dilakukan untuk mengubah data awal menjadi representasi fitur yang lebih stabil, informatif, dan sesuai untuk model klasifikasi kredit. Tahap ini sangat penting karena kualitas fitur memiliki pengaruh langsung terhadap akurasi model, stabilitas prediksi, dan kemudahan interpretasi hasil. Dalam penelitian ini, tahapan pra-pemrosesan mencakup pembersihan data, transformasi fitur, seleksi fitur, dan penanganan ketidakseimbangan kelas.

\subsubsection{Pembersihan Data dan Rekayasa Fitur}
Pada tahap awal, dilakukan pemeriksaan konsistensi tipe data, kelengkapan nilai, dan validitas rentang numerik. Nilai yang hilang atau tidak konsisten ditangani sesuai sifat variabelnya, baik melalui imputasi sederhana, pengelompokan ulang kategori, maupun penghapusan observasi yang tidak layak digunakan. Setelah itu, fitur-fitur turunan yang relevan dibentuk, seperti rasio angsuran terhadap pendapatan dan indikator beban utang, untuk memperkuat sinyal prediksi terhadap risiko gagal bayar.

\subsubsection{Transformasi Weight of Evidence (WoE)}
Seluruh fitur yang relevan dikonversi menggunakan teknik \textit{Weight of Evidence} (WoE) untuk menangkap hubungan nonlinier dan meningkatkan stabilitas prediksi[cite: 43]. Transformasi ini dilakukan melalui proses pengelompokan nilai ke dalam beberapa \textit{bin}, kemudian menghitung rasio logaritmik antara proporsi debitur baik dan debitur gagal bayar pada setiap kelompok. Secara matematis, WoE dihitung sebagai berikut[cite: 51]:
\begin{equation}
    WoE_{i} = \ln \left( \frac{N_{good,i} / N_{good,total}}{N_{bad,i} / N_{bad,total}} \right)
\end{equation}
[cite: 53]

Penggunaan WoE dalam penelitian ini bertujuan untuk menghasilkan fitur yang lebih mudah diinterpretasikan dan lebih stabil terhadap perubahan distribusi data. Selain itu, transformasi ini memudahkan analisis drift pada tahap monitoring karena perubahan proporsi kategori dapat diterjemahkan ke dalam perubahan nilai WoE secara kuantitatif.

\subsubsection{Information Value (IV) dan Seleksi Fitur}
Setelah nilai WoE diperoleh, daya prediksi setiap variabel diukur menggunakan metrik \textit{Information Value} (IV) untuk mengurangi kompleksitas model[cite: 65]. Variabel dengan nilai IV antara 0.1--0.5 dikategorikan sebagai prediktor menengah hingga kuat[cite: 70]. Variabel dengan daya prediksi sangat rendah dieliminasi agar model tidak dipengaruhi oleh fitur yang tidak informatif, sedangkan variabel yang terlalu dominan dievaluasi kembali untuk menghindari potensi \textit{data leakage}. Melalui tahapan ini, himpunan fitur akhir diharapkan mampu menjaga keseimbangan antara akurasi, stabilitas, dan interpretabilitas model.

\subsubsection{Penanganan Ketidakseimbangan Kelas}
Karena proporsi nasabah gagal bayar lebih kecil dibandingkan nasabah lancar, dilakukan penanganan \textit{class imbalance} agar model tidak bias terhadap kelas mayoritas. Dalam penelitian ini, pendekatan yang digunakan adalah \textit{Synthetic Minority Oversampling Technique} (SMOTE) pada data latih. SMOTE diterapkan setelah proses pembagian data agar informasi sintetis dari kelas minoritas tidak bocor ke data validasi maupun data uji. Dengan pendekatan ini, model diharapkan memiliki sensitivitas yang lebih baik dalam mendeteksi kasus berisiko tinggi tanpa mengorbankan kemampuan generalisasi.

\subsection{Pengembangan Model Risiko Kredit}
Tahap pengembangan model bertujuan membangun model klasifikasi yang mampu memprediksi probabilitas gagal bayar secara akurat dan konsisten. Penelitian ini menggunakan algoritma XGBoost sebagai model utama karena algoritma tersebut memiliki kemampuan yang baik dalam menangani hubungan nonlinier, interaksi antarfitur, dan struktur data tabular yang kompleks. Selain itu, XGBoost dikenal efisien secara komputasional dan sering digunakan dalam berbagai studi risiko kredit karena performanya yang kompetitif.

\subsubsection{Pelatihan Model}
Pelatihan model dilakukan menggunakan data latih yang telah melalui proses pra-pemrosesan. Pada tahap ini, model mempelajari pola hubungan antara fitur-fitur nasabah dan label probabilitas gagal bayar. Proses pelatihan dilakukan secara iteratif dengan mengatur kombinasi parameter tertentu, seperti kedalaman pohon, \textit{learning rate}, jumlah estimator, dan \textit{subsample ratio}. Setiap konfigurasi yang diuji dicatat agar dapat dibandingkan secara sistematis pada tahap evaluasi.

\subsubsection{Optimasi Hyperparameter}
Untuk memperoleh konfigurasi model yang optimal, penelitian ini menggunakan pendekatan \textit{RandomizedSearchCV}[cite: 15]. Metode ini dipilih karena lebih efisien dibandingkan pencarian menyeluruh pada ruang parameter yang luas. Kombinasi parameter terbaik dipilih berdasarkan performa pada data validasi, terutama dari sisi diskriminasi model terhadap kelas gagal bayar dan kelas tidak gagal bayar. Dengan strategi ini, model yang diperoleh tidak hanya baik pada data pelatihan, tetapi juga memiliki kemampuan generalisasi yang lebih kuat.

\subsubsection{Kriteria Pemilihan Model}
Model terbaik dipilih berdasarkan kombinasi beberapa metrik evaluasi, bukan hanya satu ukuran tunggal. Pendekatan ini digunakan karena penilaian risiko kredit menuntut keseimbangan antara kemampuan mendeteksi debitur berisiko dan menjaga agar keputusan penolakan kredit tidak berlebihan. Oleh sebab itu, kriteria pemilihan model mempertimbangkan metrik diskriminasi, metrik klasifikasi, dan kemudahan interpretasi agar model yang dihasilkan tetap relevan baik secara teknis maupun operasional.

\subsection{Validasi dan Evaluasi Model}
Evaluasi model dilakukan untuk menilai sejauh mana model mampu menghasilkan prediksi yang akurat, stabil, dan dapat digunakan dalam pengambilan keputusan kredit. Pada penelitian ini, evaluasi dilakukan terhadap performa prediktif model sekaligus terhadap kualitas data sintetis yang digunakan selama proses pengembangan.

\subsubsection{Evaluasi Performa Prediktif}
Kinerja model diukur menggunakan beberapa metrik, yaitu \textit{Area Under the Curve} (AUC), \textit{precision}, \textit{recall}, \textit{F1-score}, dan \textit{confusion matrix}. AUC digunakan untuk menilai kemampuan model dalam membedakan debitur berisiko dan tidak berisiko pada berbagai ambang keputusan. \textit{Precision} menunjukkan ketepatan model ketika menandai debitur sebagai berisiko, sedangkan \textit{recall} menunjukkan kemampuan model menangkap sebanyak mungkin kasus gagal bayar. \textit{F1-score} digunakan untuk menyeimbangkan kedua aspek tersebut, sementara \textit{confusion matrix} memberikan gambaran rinci terhadap kesalahan klasifikasi yang terjadi.

Selain metrik umum, penentuan ambang klasifikasi juga dipertimbangkan sebagai bagian penting dari evaluasi. Ambang keputusan tidak selalu ditetapkan pada nilai default 0.5, melainkan dapat disesuaikan berdasarkan toleransi risiko lembaga keuangan. Dalam konteks bisnis, kesalahan menerima debitur berisiko tinggi dapat berdampak pada peningkatan kredit bermasalah, sedangkan kesalahan menolak debitur yang sebenarnya sehat dapat mengurangi peluang bisnis. Oleh sebab itu, pemilihan ambang keputusan dilakukan dengan memperhatikan keseimbangan antara risiko dan peluang.

\subsubsection{Validasi Data Sintetis (Quality Gate)}
Sistem ini mengimplementasikan \textit{quality gate} yang ketat sebelum model dideploy ke produksi[cite: 94]. Validasi ini bertujuan memastikan bahwa data sintetis tidak hanya menyerupai data asli secara statistik, tetapi juga cukup aman untuk digunakan tanpa membocorkan informasi sensitif. Komponen validasi utama terdiri atas:
\begin{itemize}
    \item \textbf{Validasi Utilitas}: menggunakan metode \textit{Train Synthetic Test Real} (TSTR) dengan ambang batas $\Delta AUC \leq 10\%$[cite: 97, 101]. Pengujian ini menilai apakah model yang dilatih pada data sintetis masih mampu bekerja dengan baik ketika diuji pada data yang merepresentasikan kondisi riil.
    \item \textbf{Uji Distribusi}: menggunakan \textit{Inverted KS Statistic} dengan ambang batas \textit{p-value} 0.05[cite: 105, 111]. Pengujian ini bertujuan melihat kesamaan distribusi antarvariabel penting antara data sintetis dan data acuan.
    \item \textbf{Verifikasi Privasi}: menggunakan \textit{Distance to Closest Record} (DCR) untuk memastikan tidak terjadi \textit{memorization} terhadap data asli[cite: 112, 113]. Pengujian ini penting agar data sintetis yang digunakan tetap aman dari risiko reproduksi identitas individual.
\end{itemize}

Apabila salah satu kriteria pada \textit{quality gate} tidak terpenuhi, maka dataset sintetis dinyatakan belum layak digunakan dalam pipeline berikutnya. Konsekuensinya, proses pembangkitan data perlu diulang atau parameter generator perlu disesuaikan hingga data memenuhi ambang utilitas dan privasi yang telah ditetapkan. Mekanisme ini memastikan bahwa model yang dibangun di atas data sintetis tetap memiliki kredibilitas ilmiah dan operasional.

\subsection{Implementasi MLOps}
Setelah model terbaik diperoleh, tahap berikutnya adalah mengintegrasikan model ke dalam pipeline MLOps agar dapat dioperasikan secara berkelanjutan. Implementasi MLOps dalam penelitian ini mencakup orkestrasi workflow, pelacakan eksperimen, registrasi model, deployment layanan inferensi, pemantauan drift, dan retraining otomatis. Seluruh proses tersebut disusun untuk meminimalkan pekerjaan manual dan mempercepat respons sistem terhadap perubahan data.

\subsubsection{Orkestrasi Workflow dan Pelacakan Eksperimen}
\textit{Workflow} utama diatur menggunakan n8n, yang berfungsi memicu dan menghubungkan setiap tahapan dalam \textit{pipeline}. Sebagai contoh, n8n dapat menjalankan proses pelatihan saat dataset baru tersedia, melanjutkan ke proses evaluasi, lalu mendaftarkan model terbaik ke registry bila memenuhi kriteria performa. Sementara itu, MLflow digunakan untuk menyimpan parameter eksperimen, metrik hasil pelatihan, dan artefak model. Dengan demikian, setiap eksperimen dapat ditelusuri kembali sehingga proses pengembangan model menjadi lebih transparan dan reprodusibel.


Selain \textit{pipeline} \textit{retraining}, \textit{workflow} n8n juga mengintegrasikan jalur \textbf{ingesti data nasabah berbasis email} secara otomatis sebagai \textit{pipeline} kedua yang terpisah. Setiap jam, n8n memonitor \textit{inbox} Gmail untuk email bersubjek \textit{"Credit Application"} yang menyertakan lampiran berformat PDF. Formulir yang diterima secara otomatis diunduh, diteruskan ke skrip parse\_pdf.py untuk proses ekstraksi data, kemudian disimpan ke database nasabah melalui \textit{endpoint} POST /api/nasabah. Dengan demikian, proses \textit{onboarding} nasabah baru dapat dilakukan sepenuhnya tanpa intervensi manual --- formulir kredit yang dikirimkan melalui email langsung masuk ke database dan siap untuk diprediksi oleh model.

\begin{table}[H]
\centering
\caption{Dua \textit{Pipeline} Utama n8n dalam Sistem MLOps}
\label{tab:n8n-pipelines}
\begin{tabularx}{\textwidth}{|p{3.5cm}|X|p{3.5cm}|}
\hline
\textbf{Pipeline} & \textbf{Alur} & \textbf{Trigger} \\ \hline
\textit{Retraining Pipeline} & Schedule → Jalankan trainingCredit.py → Cek AUC $>$ 0.75 → Promosi ke Production MLflow → \textit{Reload} model & Terjadwal (7 hari sekali) \\ \hline
\textit{Ingestion Pipeline} & Gmail Trigger → Unduh lampiran PDF → parse\_pdf.py → Validasi JSON → POST /api/nasabah & Setiap jam (polling Gmail) \\ \hline
\end{tabularx}
\end{table}

Pelacakan eksperimen dilakukan menggunakan \textbf{MLflow}, yang mencatat parameter \textit{hyperparameter}, metrik evaluasi (AUC, F1-Score), dan artefak model setiap kali proses pelatihan dijalankan. Selain pencatatan eksperimen, sistem juga menggunakan \textbf{MLflow Tracing} (@mlflow.trace) untuk mencatat setiap permintaan inferensi ke \textit{endpoint} /predict beserta versi model yang digunakan. Kemampuan \textit{audit trail} per-prediksi ini mendukung kebutuhan kepatuhan regulasi, termasuk ketentuan transparansi model sesuai POJK Nomor 29 Tahun 2024.

\subsubsection{Deployment Model dan Layanan Inferensi}
Model yang lolos evaluasi disimpan sebagai model kandidat dan dipublikasikan sebagai layanan inferensi menggunakan FastAPI. Layanan ini menerima input atribut nasabah, melakukan transformasi yang diperlukan, kemudian menghasilkan skor risiko atau probabilitas gagal bayar. Agar konsisten di berbagai lingkungan, seluruh komponen deployment dikemas menggunakan Docker. Pendekatan kontainerisasi ini mempermudah proses distribusi, pengujian, dan migrasi sistem ke server pengembangan maupun produksi.


\subsubsection{Mekanisme Persetujuan Kredit Berbasis Email (\textit{Human-in-the-Loop})}

Sistem tidak memberikan keputusan kredit secara otomatis. Setiap hasil prediksi model harus dikonfirmasi secara manual oleh analis kredit atau pejabat pemutus melalui mekanisme \textit{Human-in-the-Loop} (HITL) berbasis email. Mekanisme ini diimplementasikan melalui alur berikut:

\begin{enumerate}
    \item Setelah model menghasilkan prediksi, analis kredit memicu \textit{endpoint} POST /api/\allowbreak nasabah/\allowbreak\{id\/\allowbreak send-confirmation} dengan menyertakan alamat email reviewer.
    \item Sistem menyusun \textbf{email HTML} yang berisi ringkasan risiko nasabah---mencakup \textit{bureau score}, pendapatan bulanan, jumlah pinjaman, \textit{risk grade}, dan probabilitas \textit{default}---beserta dua tombol interaktif: \textbf{Setuju} dan \textbf{Tolak}.
    \item Email dikirimkan kepada reviewer melalui \textbf{n8n webhook} (/webhook/send-confirmation) yang meneruskan pesan via server SMTP Gmail.
    \item Setiap pengiriman email dicatat secara otomatis ke tabel sent\_emails sebagai \textit{audit trail}.
    \item Klik pada tombol Setuju memanggil GET /api/nasabah/\{id\/approve}, sedangkan klik Tolak memanggil GET /api/nasabah/\{id\/reject}. Kedua \textit{endpoint} ini memperbarui kolom approval\_status di database menjadi "Accepted" atau "Rejected" dan menampilkan halaman konfirmasi kepada reviewer.
    \item Apabila email konfirmasi dikirim ulang untuk nasabah yang sama, sistem secara otomatis mereset approval\_status kembali ke "Pending" agar siklus peninjauan ulang dimulai dari awal.
\end{enumerate}

\begin{figure}[H]
\centering
\begin{tikzpicture}[node distance=1.0cm and 2.0cm]
  \node[brstep, align=center] (pd)     {Prediksi PD\\dan risk\_grade selesai};
  \node[brstep, align=center, below=of pd] (email)  {Kirim email konfirmasi\\via n8n (SMTP)};
  \node[brdecision, align=center, below=of email] (decide) {Keputusan\\Analis?};
  \node[brok, align=center, right=3cm of decide] (acc)  {approval\_status\\= \textbf{Accepted}};
  \node[brerr, align=center, left=3cm of decide] (rej)  {approval\_status\\= \textbf{Rejected}};
  \node[brstep, align=center, above=1.2cm of email] (reset) {Reset ke\\approval\_status\\= \textbf{Pending}};

  \draw[brarrow] (pd)     -- (email);
  \draw[brarrow] (email)  -- (decide);
  \draw[brarrow] (decide) -- node[above,font=\footnotesize]{Setuju} (acc);
  \draw[brarrow] (decide) -- node[above,font=\footnotesize]{Tolak} (rej);
  \draw[brarrow, dashed, bend left=50] (email.east) to
    node[right, align=center, font=\footnotesize]{Kirim ulang\\email} (reset.east);
  \draw[brarrow, dashed] (reset) -- (email);
\end{tikzpicture}
\caption{Siklus approval\_status --- keputusan kredit berbasis konfirmasi email manusia.}
\label{fig:hitl-flow}
\end{figure}

Mekanisme ini memastikan bahwa model \textit{machine learning} hanya berperan sebagai \textbf{alat bantu keputusan}, bukan pengganti penilaian manusia. Seluruh keputusan final tetap berada di tangan analis kredit yang berwenang.

\subsubsection{Sistem Grading Risiko Kredit}
Output layanan inferensi tidak hanya berupa nilai probabilitas gagal bayar, tetapi juga diterjemahkan ke dalam sistem grading agar lebih mudah digunakan dalam pengambilan keputusan operasional. Sistem grading yang digunakan terdiri atas Grade A, B, C, D, dan Undefined. Grade A menunjukkan kemungkinan gagal bayar yang rendah sehingga calon peminjam berada pada kategori risiko paling baik. Grade B dan C menunjukkan tingkat risiko menengah yang masih dapat dipertimbangkan dengan analisis tambahan sesuai kebijakan kredit. Grade D menunjukkan kemungkinan gagal bayar yang tinggi sehingga pengajuan perlu diperlakukan sebagai risiko tinggi atau berpotensi ditolak.

Kategori Undefined digunakan ketika sistem tidak dapat langsung memberikan keputusan berbasis model karena terdapat kondisi khusus yang membutuhkan intervensi analis, pengecekan manual, atau penanggulangan tambahan. Kondisi tersebut dapat terjadi ketika karakteristik pengajuan tidak cukup aman untuk dievaluasi hanya dengan probabilitas default, misalnya usia peminjam terlalu tinggi, tenor pinjaman terlalu panjang, atau usia saat pelunasan melewati batas kebijakan internal. Pada kasus seperti ini, hasil model perlu dilengkapi dengan aturan bisnis, seperti rekomendasi pengurangan tenor, pertimbangan asuransi jiwa kredit, atau eskalasi ke komite kredit.

Secara metodologis, pemisahan antara grade A--D dan Undefined penting karena model \textit{machine learning} hanya mempelajari pola statistik dari data, sedangkan keputusan kredit juga harus mempertimbangkan batasan kebijakan, risiko operasional, dan aspek kepatuhan. Dengan demikian, sistem grading berfungsi sebagai lapisan translasi dari skor prediktif menjadi keputusan yang lebih mudah dipahami, sementara kategori Undefined berfungsi sebagai mekanisme pengaman agar sistem tidak memberikan persetujuan otomatis pada kasus yang memerlukan penilaian manusia.


Penerapan batas usia pelunasan dilakukan melalui sistem logika hibrida yang menghitung usia nasabah pada saat jatuh tempo pinjaman. Formulasi yang digunakan adalah sebagai berikut:

\begin{equation}
    \text{Usia Pelunasan} = \text{Usia Saat Pengajuan} + \frac{\text{Tenor (Bulan)}}{12}
    \label{eq:usia-pelunasan}
\end{equation}

Apabila nilai usia pelunasan pada Persamaan~\ref{eq:usia-pelunasan} melebihi 80 tahun, sistem memberikan status Grade Undefined dan menghitung rekomendasi tenor maksimal yang diperbolehkan secara dinamis:

\begin{equation}
    \text{Tenor Maksimal (Bulan)} = (80 - \text{Usia Saat Pengajuan}) \times 12
    \label{eq:tenor-maks}
\end{equation}

Berdasarkan kedua persamaan tersebut, sistem membagi kondisi nasabah ke dalam \textbf{tiga zona keputusan}:

\begin{enumerate}
    \item \textbf{Zona Lolos} --- Usia pelunasan $\leq$ 80 tahun. Data diteruskan ke model XGBoost secara normal untuk memperoleh \textit{risk grade} A, B, C, atau D.
    \item \textbf{Zona Mitigasi Dinamis} --- Usia nasabah saat pengajuan masih di bawah 80 tahun, tetapi tenor yang diminta menyebabkan usia pelunasan melewati 80 tahun. Sistem memberikan Grade Undefined dan menghitung rekomendasi tenor maksimal menggunakan Persamaan~\ref{eq:tenor-maks}. Jika nasabah memiliki asuransi jiwa kredit aktif yang nilainya mencukupi jumlah pinjaman, sistem juga dapat merekomendasikan dispensasi melalui komite kredit.
    \item \textbf{Zona \textit{Joint-Borrower}} --- Usia nasabah saat pengajuan telah mencapai 80 tahun atau lebih. Pada kondisi ini, pengurangan tenor tidak lagi memadai karena sisa batas usia bernilai nol atau negatif. Sistem merekomendasikan alternatif pengajuan berbasis \textit{joint-income} atau \textit{co-borrower} bersama anggota keluarga inti yang berusia di bawah 60 tahun.
\end{enumerate}

Pemisahan tiga zona ini penting karena model \textit{machine learning} hanya mempelajari korelasi statistik antarfitur dan tidak secara eksplisit memahami batas kelayakan usia pelunasan yang diatur oleh kebijakan manajemen risiko institusi. Dengan demikian, logika bisnis pada lapisan API berfungsi sebagai mekanisme pengaman yang melengkapi prediksi model.

\subsubsection{Monitoring Drift dan Retraining Otomatis}
Setelah model aktif digunakan, sistem melakukan pemantauan terhadap perubahan distribusi fitur dan perubahan hubungan antara fitur dengan target. Pemantauan \textit{data drift} dilakukan dengan membandingkan distribusi data baru terhadap distribusi referensi, sedangkan pemantauan \textit{concept drift} dilakukan dengan mengamati penurunan metrik kinerja model dari waktu ke waktu. Jika indikator \textit{drift} melewati ambang batas yang telah ditentukan, n8n akan memicu ulang proses pelatihan model menggunakan data terbaru. Dengan mekanisme ini, model dapat diperbarui secara lebih responsif tanpa menunggu evaluasi manual yang berpotensi terlambat.


Pemantauan \textit{drift} pada sistem ini diimplementasikan menggunakan pendekatan \textbf{dual-trigger} dalam modul drift\_monitor.py. Sistem tidak bergantung pada satu mekanisme tunggal, melainkan menggabungkan dua jalur pemicuan yang saling melengkapi:

\begin{enumerate}
    \item \textbf{Trigger Berbasis Jadwal (\textit{Schedule-Based Trigger})} --- n8n menjalankan \textit{workflow} pemantauan secara terjadwal setiap \textbf{7 hari sekali}. Pada setiap eksekusi, modul membandingkan distribusi data inferensi terkini dengan distribusi data referensi menggunakan \textit{Kolmogorov-Smirnov (KS) Statistic}\footnote{KS Statistic adalah uji non-parametrik yang mengukur jarak maksimum antara dua fungsi distribusi kumulatif (CDF). Penggunaannya tidak memerlukan asumsi distribusi tertentu, sehingga cocok untuk fitur tabular dengan distribusi beragam.}. Nilai KS-statistic dihitung per fitur; apabila KS-statistic pada fitur kritis (seperti monthly\_income atau loan\_amount) melampaui ambang batas $\alpha = 0.05$ berdasarkan uji signifikansi statistik, sistem mengklasifikasikan kondisi sebagai \textit{drift terdeteksi}.

    \item \textbf{Trigger Berbasis Event (\textit{Event-Based Trigger})} --- Setiap kali data nasabah baru dikirim melalui \textit{endpoint} POST /api/nasabah, sistem secara otomatis menghitung jumlah kumulatif data baru yang tersimpan sejak model terakhir dilatih. Apabila jumlah data baru melampaui ambang batas konfigurasi (\textit{ingestion threshold}), \textit{workflow} n8n memicu proses evaluasi dan retraining secara langsung tanpa menunggu jadwal mingguan.
\end{enumerate}

Kombinasi kedua trigger ini memastikan sistem responsif terhadap perubahan mendadak (ditangani oleh event trigger) sekaligus melakukan pemantauan rutin terjadwal (ditangani oleh schedule trigger). Pada lingkungan produksi, mekanisme ini mencegah model berjalan terlalu lama pada distribusi data yang telah bergeser signifikan.

Agar deteksi \textit{drift} dapat beroperasi secara \textit{closed-loop}, setiap permintaan prediksi ke \textit{endpoint} /predict secara otomatis dicatat ke file n8n\_data/inference\_logs.jsonl. File \textit{log} ini berisi seluruh fitur input nasabah, probabilitas \textit{default} yang dihasilkan model, dan \textit{timestamp} prediksi. Modul drift\_monitor.py kemudian membandingkan distribusi data di file \textit{log} ini dengan distribusi referensi dari data pelatihan. Tanpa mekanisme pencatatan inferensi ini, deteksi \textit{drift} tidak memiliki data produksi yang representatif untuk dianalisis.

\subsection{Interpretasi Model (XAI) dan Kepatuhan Regulasi}
Untuk menjaga transparansi dan mendukung kebutuhan audit, sistem menggunakan metode SHAP (\textit{SHapley Additive exPlanations})[cite: 149]. SHAP digunakan untuk menjelaskan kontribusi masing-masing fitur terhadap skor risiko yang dihasilkan model. Penjelasan ini penting karena model risiko kredit tidak hanya dituntut akurat, tetapi juga harus dapat memberikan alasan yang dapat dipahami oleh analis, auditor, maupun pihak regulator.

Dalam konteks implementasi, penjelasan SHAP dapat ditampilkan dalam bentuk kontribusi global maupun lokal. Penjelasan global digunakan untuk mengidentifikasi variabel paling berpengaruh terhadap keputusan model secara keseluruhan, sedangkan penjelasan lokal digunakan untuk menjelaskan alasan di balik prediksi pada satu nasabah tertentu. Pendekatan ini relevan untuk memenuhi kebutuhan kepatuhan, termasuk mandat POJK Nomor 29 Tahun 2024 yang menekankan pentingnya transparansi pada keputusan skor kredit[cite: 156, 157]. Dengan adanya kemampuan interpretasi tersebut, sistem yang dikembangkan menjadi lebih akuntabel dan lebih siap diadopsi dalam lingkungan operasional yang diatur secara ketat.

\subsection{Ringkasan Bab}
Bab ini telah menjelaskan metodologi penelitian yang digunakan untuk membangun sistem \textit{Credit Risk Scoring} berbasis MLOps secara menyeluruh. Tahapan penelitian mencakup simulasi data, perancangan database, pra-pemrosesan, pengembangan model, validasi kualitas data sintetis, implementasi workflow MLOps, mekanisme persetujuan berbasis manusia, serta interpretasi model untuk mendukung kepatuhan regulasi. Seluruh rangkaian metodologi tersebut menjadi dasar untuk pelaksanaan pengujian dan analisis hasil pada bab berikutnya.